\documentclass[12pt,a4paper]{article}
\usepackage[T1]{fontenc}
\usepackage[utf8]{inputenc}
\usepackage[polish]{babel}
\usepackage[a4paper,margin=2.5cm]{geometry}
\usepackage{amsmath,amssymb}
\usepackage{listings}
\usepackage{pgfplots}
\pgfplotsset{width=10cm,compat=1.9}
\usepackage{pdfcomment}
\usepackage{svg}
\usepackage{float}


% Externalization requires -shell-escape; disable to avoid build errors.
%\usepgfplotslibrary{external}
%\tikzexternalize

\linespread{0.96} % lekko ciaśniej (bez psucia czytelności)

\title{\textbf{Sprawozdanie 1}\\Wybrane Zagadnienia Algebry}
\author{Jakub Kowal\\279706}
\date{}

%https://tex.stackexchange.com/questions/212793/how-can-i-typeset-julia-code-with-the-listings-package

\lstdefinelanguage{Julia}%
  {morekeywords={abstract,break,case,catch,const,continue,do,else,elseif,%
      end,export,false,for,function,immutable,import,importall,if,in,%
      macro,module,otherwise,quote,return,switch,true,try,type,typealias,%
      using,while},%
   sensitive=true,%
   alsoother={$},%
   morecomment=[l]\#,%
   morecomment=[n]{\#=}{=\#},%
   morestring=[s]{"}{"},%
   morestring=[m]{'}{'},%
}[keywords,comments,strings]%

\begin{document}
\maketitle

Użyte parametry do testów:
\[
a=2,\ b=7,\ c=9,\ d=7,\ e=0,\ f=6.
\]

\section*{Zadanie 1 z listy 2}

\subsection*{a)}
Tworzymy strukturę $Wielomian$, która przechowuje jednomiany w strukturze \pdftooltip{\textit{Dict*}}{Struktura \textit{słownik} podobna do \textit{map} znanych z innych języków programowania} w postaci wektorów z wykładnikami potęg zmiennych oraz współczynnika przy nich. Dodatkowo przechowuje również liczbę tych zmiennych.

\subsection*{b)}
\subsubsection*{Standardowy Lex}
Standardowy Lex porównuje jednomiany na podstawie wektora wykładników. Idąc ustaloną kolejnością, jeżeli jeden z wykładników będzie mniejszy niż wykładnik tej samej zmiennej w drugim jednomianie, to wiemy, że drugi $<_{\mathrm{lex}}$ pierwszy. Przykład: 

\[
x^2y\ \text{vs}\ xy^2:\quad
\alpha=(2,1),\ \beta=(1,2),\ \text{pierwsza różnica przy }i=1:\ 2>1
\]
czyli \(xy^2 <_{\mathrm{lex}} x^2y\).

\[
x^2\ \text{vs}\ x^2y:\quad
\alpha=(2,0),\ \beta=(2,1),\ \text{pierwsza różnica przy }i=2:\ 0<1
\]
czyli \(x^2 <_{\mathrm{lex}} x^2y\).

\subsubsection*{Permutowany Lex}
Działa na takiej zasadzie jak Standardowy Lex, ale permutujemy kolejność zmiennych, czyli "priorytet" wykładników.

\subsubsection*{GradedLex}
GradedLex działa na takiej samej zasadzie co Standardowy Lex, ale w pierwszej kolejności bierze pod uwagę sumę wykładników. 

\subsection*{c)}

\subsection*{Opis funkcji \texttt{PolynomialReduce}}
Funkcja \texttt{PolynomialReduce} realizuje redukcje wielomianu $f$ względem zbioru wielomianow $G$ w ustalonym porzadku jednomianow. Zwraca wektory wielomianow $\alpha_1,\dots,\alpha_m$ oraz reszte $r$ takie, ze
\[
f = \sum_{i=1}^{m} \alpha_i g_i + r,
\]
przy czym żadny jednomian z $r$ nie jest podzielny przez wyraz wiodacy żadnego $g_i$ (wzgladem zadanego porzadku). Algorytm działa iteracyjnie na kopii $p$ wielomianu $f$: w każdym kroku bierze wyraz wiodący $(x^{e_p}, c_p)$, próuje znalezc $g_i$ z $G$ o wyrazie wiodącym $(x^{e_g}, c_g)$ dzielącym $x^{e_p}$, a następnie odejmuje odpowiedni wielomian $t\cdot g_i$, gdzie
\[
t = \frac{c_p}{c_g} x^{e_p - e_g}.
\]
Jesli taka redukcja jest możliwa, aktualizuje $p$ i $\alpha_i$; w przeciwnym razie przenosi wyraz wiodący do reszty $r$.

Kluczowy fragment implementacji (pętla redukcji):
\begin{lstlisting}[language=Julia]
while !isempty(p.coeffs)
    e_p, c_p = leading_term(p, order)
    reduced = false

    for (i, g) in enumerate(G)
        isempty(g.coeffs) && continue
        e_g, c_g = leading_term(g, order)
        if monomial_divides(e_g, e_p)
            exp_diff = sub_exp(e_p, e_g)
            coeff_mult = c_p / c_g
            t = monomial(coeff_mult, exp_diff, f.nvars)
            alphas[i] = alphas[i] + t
            p = p - t * g
            reduced = true
            break
        end
    end

    if !reduced
        lt = monomial(c_p, e_p, f.nvars)
        r = r + lt
        p = p - lt
    end
end
\end{lstlisting}

\subsubsection*{d)}
Zadanie 37 z ćwiczeń:
\begin{center}
    $f = x^{3} - x^{2}y - x^{2}z$\\
$g_{1} = x^{2}y - z$\\
$g_{2} = xy - 1$
\end{center}
Porządek GradedLex (z $<_{\mathrm{lex}}$ y $<_{\mathrm{lex}}$ x)\\
\begin{enumerate}
    \item Redukcja f przez ($g_{1},g_{2}$)
    \item Redukcja f przez ($g_{2},g_{1}$)
    \item Czy $r_{1} - r_{2} \in <g_{1},g_{2}>$?
\end{enumerate}

Wyniki:
\begin{enumerate}
    \item $f=-1 * g_{1} + 0*g_{2} + r_{1}$\\$r_{1} = x^{3} - x^{2}z - z$
    \item $f=-x * g_{2} + 0*g_{1} + r_{2}$\\$r_{2} = x^{3} - x^{2}z - x$
    \item $r_{1} - r_{2} \notin <g_{1},g_{2}>$
\end{enumerate}

Różnica wynika z tego, że \textit{PolynomialReduce} jest algorytmem zależnym od kolejności dzielników: w każdej iteracji bierze pierwszy wielomian z listy, którego wyraz wiodący dzieli wyraz wiodący reszty. Po zamianie $g_{1}$ i $g_{2}$ zmienia się, który reduktor „złapie” jako pierwszy, więc ścieżka redukcji i współczynniki w alphas są inne, a także sama reszta ($r_{1}$ i $r_{2}$).

\subsubsection*{e)}
Dla danych $a = 2,b= 7,c = 9,d = 7,e = 0,f= 6$ mamy
\[
h = x^{2}y^{7} + y^{9}z^{7} + z^{6}
\]
oraz przykład
\[
G = \{x^{2}y^{7} - z^{6},\; y^{9}z^{7} - x^{2}y^{7}\}.
\]
Wyniki redukcji (różne porządki leksykograficzne):
\begin{enumerate}
    \item Lex $z<_{\mathrm{lex}}y<_{\mathrm{lex}}x$:\quad $r = y^{9}z^{7} + 2z^{6}$
    \item Lex $x<_{\mathrm{lex}}z<_{\mathrm{lex}}y$:\quad $r = 3z^{6}$
    \item Lex $y<_{\mathrm{lex}}x<_{\mathrm{lex}}z$:\quad $r = x^{2}y^{16}z + 2x^{2}y^{7}$
\end{enumerate}

Różnice wynikają z tego, że porządek leksykograficzny zmienia wyraz wiodący zarówno w $h$, jak i w wielomianach z $G$. To wpływa na to, który dzielnik jest wybierany w algorytmie oraz jak przebiega ścieżka redukcji.


\section*{Zadanie 2 z listy 2}

\def\a{0}
\def\b{3.2}
\def\m{3.8}
\def\p{0}

\begin{center}
    % $r=\a \cos(\frac{2}{\theta}) + \b \cos(\m * \theta + \p)$
    $r=\b \sin(\m * \theta)$
\end{center}

\begin{center}
    \begin{tikzpicture}
    \begin{axis}[
        axis equal,
        grid=both,
        samples=1200,
        domain=0:10*pi,
    ]
    \addplot[color=red, parametric] (
        {(\b*sin(deg(\m*x)))*cos(deg(x))},
        {(\b*sin(deg(\m*x)))*sin(deg(x))}
    );
    \end{axis}
    \end{tikzpicture}
\end{center}


Dla porównania godło Japonii:

\begin{center}
   \begin{figure}[H]
        \caption{Godło Japonii}
        \hspace*{0.20\textwidth}
        \includesvg[width=250pt]{Imperial_Seal_of_Japan.svg}
        \label{rys:godlo}
    \end{figure} 
\end{center}



\section*{Wykorzystanie Sztucznej Inteligencji}
Wykorzystałem sztuczną inteligencję w ramach wsparcia redakcji opisów implementacji oraz usprawnienia przeszukiwania dokumentacji języka Julia
\end{document}
